\hspace*{1.6cm}เมื่อประเทศไทยกำลังเข้าสู่สังคมผู้สูงอายุ ความเหงาและความโดดเดี่ยวทางสังคมในหมู่ผู้สูงอายุได้กลายเป็นปัญหาสำคัญ โซลูชันการติดตามดูแลที่มีอยู่มักขาดการโต้ตอบหรือมีความซับซ้อนเกินไปสำหรับผู้ใช้สูงวัย โครงงานนี้นำเสนอ "Buddy" ระบบเพื่อนคู่ใจอัจฉริยะที่ออกแบบมาเพื่อยกระดับคุณภาพชีวิตของผู้สูงอายุผ่านการสื่อสารที่เป็นธรรมชาติและการติดตามสภาวะอารมณ์ ระบบประกอบด้วยแอปพลิเคชันบนแท็บเล็ตสำหรับผู้สูงอายุและแอปพลิเคชันมือถือสำหรับผู้ดูแล เพื่อช่วยให้การปฏิสัมพันธ์เป็นไปอย่างราบรื่น สถาปัตยกรรมทางเทคนิคใช้กระบวนการ Speech-to-Speech (STS) ที่ครบวงจร โดยใช้ Faster Whisper เพื่อการรู้จำเสียงภาษาไทยที่มีความแม่นยำสูง และโมเดลภาษาขนาดใหญ่ (Large Language Model) แบบประมวลผลภายในเครื่อง (GEMMA 3:4b) เพื่อสร้างบทสนทนาที่เข้าใจความรู้สึกและมีบุคลิกภาพที่เหมาะสม นอกจากนี้ยังมีโมเดล GEMMA 3:1b แยกต่างหากสำหรับวิเคราะห์บันทึกการสนทนาเพื่อติดตามแนวโน้มทางอารมณ์ ในขณะที่ ZegoCloud SDK ช่วยให้การวิดีโอคอลแบบเรียลไทม์เป็นไปอย่างเสถียร

\hspace*{1.6cm}ผลการทดลองยืนยันประสิทธิภาพของระบบในหลายมิติ โมดูลการรู้จำเสียงมีอัตราความผิดพลาดของคำ (WER) อยู่ที่ 10.50\% และอัตราความผิดพลาดของตัวอักษร (CER) ที่ 7.20\% ซึ่งดีกว่าเกณฑ์มาตรฐานของ Thonburian การจำแนกอารมณ์แสดงความน่าเชื่อถือสูงด้วยคะแนน F1-score เฉลี่ย 0.93 แม้ว่าโมเดล LLM ขนาดเล็กแบบประมวลผลภายในเครื่องจะมีคะแนนด้านความรู้ทั่วไปต่ำกว่าโมเดลบนคลาวด์ แต่กลับมีประสิทธิภาพสูงในการสนับสนุนทางอารมณ์และบทสนทนาในชีวิตประจำวัน โดยได้คะแนน Mean Opinion Score (MOS) จากมนุษย์ที่ 4.2 จาก 5 คะแนน ระบบการสื่อสารโทรคมนาคมแสดงประสิทธิภาพที่แข็งแกร่งด้วยอัตราความสำเร็จในการเชื่อมต่อ 98.5\% และมีความหน่วงเฉลี่ย 250 มิลลิวินาที ผลสำรวจผู้ใช้ระบุคะแนนความพึงพอใจอยู่ที่ 5.5 จาก 7 คะแนน ซึ่งชี้ให้เห็นว่า "Buddy" ประสบความสำเร็จในการลดอุปสรรคทางเทคโนโลยีและให้ข้อมูลเชิงลึกทางอารมณ์ที่มีค่าแก่ผู้ดูแล ทำหน้าที่เป็นสะพานเชื่อมระหว่างผู้สูงอายุและครอบครัวได้อย่างมีประสิทธิภาพ